% !TEX root = ../../ctfp-print.tex

\lettrine[lhang=0.17]{様}{々な例を学ぶことで、}圏についての真の理解が得られる。圏は
あらゆる形やサイズで存在し、予期しない場所によく現れる。非常に単純なものから始めよう。

\section{対象なし}

最も自明な圏は対象が0個、結果として射も0個のものだ。それ自体では非常に寂しい圏だが、
他の圏の文脈では重要かもしれない。例えば、すべての圏の圏（そう、そういうものがある）
において。空集合が意味をなすと思うなら、空の圏はどうだろう？

\section{単純なグラフ}

対象を矢印で繋ぐだけで圏を構築できる。任意の有向グラフから始め、単純に矢印を追加する
ことで圏にすることができると想像できる。まず、各ノードに恒等射を追加する。次に、一方
の終点が他方の始点と一致する2つの矢印（言い換えれば、\newterm{合成可能な}2つの矢印）
に対して、それらの合成として機能する新しい矢印を追加する。新しい矢印を追加するたびに、
他のすべての矢印（恒等射を除く）とその矢印自身との合成も考慮しなければならない。通常は
無限に多くの矢印で終わるが、それでも構わない。

このプロセスを別の観点から見ると、グラフの各ノードに対応する対象と、合成可能なグラフ
の辺のすべての可能な\newterm{鎖}を射として持つ圏を作っていることになる。（恒等射を長さ
ゼロの鎖の特殊ケースとして考えることもできる。）

そのような圏は、与えられたグラフによって生成される\newterm{自由圏}と呼ばれる。これは
自由構成の例であり、法則（ここでは圏の法則）を満たすために最小限の項目で拡張すること
によって与えられた構造を完成させるプロセスだ。将来、その例をさらに見ることになるだろう。

\section{順序}

今度はまったく違うものを！射がオブジェクト間の関係――以下または等しいという関係――で
ある圏を見てみよう。これが実際に圏であるかどうか確認しよう。恒等射はあるか？すべての
対象は自分自身以下または等しい：合格！合成はあるか？$a \leqslant b$かつ$b \leqslant c$
なら$a \leqslant c$：合格！合成は結合的か？合格！このような関係を持つ集合は
\newterm{前順序}と呼ばれる。だから前順序は確かに圏だ。

より強い関係もある。$a \leqslant b$かつ$b \leqslant a$なら$a$は$b$と同じでなければ
ならないという追加条件を満たすものだ。それは\newterm{半順序}と呼ばれる。

最後に、任意の2つの対象がいずれかの方向で互いに関係にあるという条件を課すことができる。
それが\newterm{線形順序}または\newterm{全順序}だ。

これらの順序集合を圏として特徴付けよう。前順序は任意の対象$a$から任意の対象$b$へ向か
う射が高々1つしかない圏だ。そのような圏の別名は「細い（thin）」だ。前順序は細い圏だ。

圏$\cat{C}$における対象$a$から対象$b$への射の集合は\newterm{ホム集合}と呼ばれ、
$\cat{C}(a, b)$（または$\mathbf{Hom}_{\cat{C}}(a, b)$）と書かれる。だから前順序のすべ
てのホム集合は空か単集合だ。これには前順序において恒等射のみを含む単集合でなければ
ならない、$a$から$a$への射の集合$\cat{C}(a, a)$が含まれる。ただし、前順序にはサイクル
があってもよい。半順序ではサイクルは禁止されている。

ソートのために、前順序、半順序、全順序を認識できることは非常に重要だ。クイックソート、
バブルソート、マージソートなどのソートアルゴリズムは全順序に対してのみ正しく機能する。
半順序はトポロジカルソートを使ってソートできる。

\section{集合としてのモノイド}

モノイドは恥ずかしいほど単純だが驚くほど強力な概念だ。基本的な算術の背後にある概念だ：
加算と乗算の両方がモノイドを形成する。モノイドはプログラミングにいたるところに現れる。
文字列、リスト、畳み込み可能なデータ構造、並行プログラミングにおけるフューチャー、
関数型リアクティブプログラミングにおけるイベントなどとして現れる。

伝統的に、モノイドは二項演算を持つ集合として定義される。この演算に求められるのは、それ
が結合的であることと、それに関して単位として振る舞う特別な要素が1つ存在することだ。

例えば、ゼロを含む自然数は加算のもとでモノイドを形成する。結合性とは：
\[(a + b) + c = a + (b + c)\]
（言い換えれば、数を加算するとき括弧を省略できる。）

単位元はゼロだ。なぜなら：
\[0 + a = a\]
および
\[a + 0 = a\]
2番目の式は冗長だ。加算は可換（$a + b = b + a$）だからだ。しかし可換性はモノイドの定
義の一部ではない。例えば、文字列の連結は可換ではないが、モノイドを形成する。ちなみに
文字列連結の単位元は空文字列で、文字列のどちらの側に付けても変化しない。

Haskellではモノイドの型クラスを定義できる――\code{mempty}と呼ばれる単位元と
\code{mappend}と呼ばれる二項演算を持つ型だ：

\src{snippet01}
2引数関数の型シグネチャ\code{m -> m -> m}は最初は奇妙に見えるかもしれないが、カリー化
について話した後は完全に意味をなす。複数の矢印を持つシグネチャは2つの基本的な方法で解
釈できる：最右の型が戻り型である複数の引数を持つ関数として；あるいは関数を返す1つの引
数（最左のもの）の関数として。後者の解釈は括弧を追加することで強調できる（矢印は右結合
なので括弧は冗長だが）：\code{m -> (m -> m)}。この解釈にはすぐに戻る。

Haskellでは\code{mempty}と\code{mappend}のモノイド的性質（つまり\code{mempty}が中立的
で\code{mappend}が結合的であるという事実）を表現する方法がないことに注意しよう。それら
が満たされることを確認するのはプログラマーの責任だ。

HaskellのクラスはC++のクラスほど押し付けがましくない。新しい型を定義するとき、事前に
そのクラスを指定する必要はない。先延ばしして、後で指定した型をあるクラスのインスタンス
として宣言することは自由だ。例として、\code{mempty}と\code{mappend}の実装を提供すること
で\code{String}をモノイドとして宣言しよう（実際にはこれは標準のPreludeでやってくれて
いる）：

\src{snippet02}
ここで、\code{String}は単に文字のリストなので、リスト連結演算子\code{(++)}を再利用した。

Haskellの構文について一言：任意の中置演算子は括弧で囲むことで2引数関数に変換できる。
2つの文字列を与えられたとき、それらの間に\code{++}を挿入することで連結できる：

\begin{snip}{haskell}
"Hello " ++ "world!"
\end{snip}
あるいは括弧で囲んだ\code{(++)}に2つの引数として渡すことで：

\begin{snip}{haskell}
(++) "Hello " "world!"
\end{snip}
関数への引数はカンマで区切られず括弧で囲まれていないことに注意しよう。（これはおそらく
Haskellを学ぶ際に最も慣れるのが難しいことだ。）

Haskellが関数の等式を表現できることを強調する価値がある：

\begin{snip}{haskell}
mappend = (++)
\end{snip}
概念的にはこれは、関数の生成する値の等式を表現することとは異なる：

\begin{snip}{haskell}
mappend s1 s2 = (++) s1 s2
\end{snip}
前者は圏$\Hask$（または終了しない計算の名前であるボトムを無視すれば$\Set$）における射
の等式に変換される。そのような等式はより簡潔なだけでなく、他の圏に一般化できることが多
い。後者は\newterm{外延的}等式と呼ばれ、任意の2つの入力文字列に対して\code{mappend}と
\code{(++)}の出力が同じという事実を述べる。引数の値はときに\newterm{点}（例えば：点$x$
での$f$の値）と呼ばれるので、これは点ごとの等式と呼ばれる。引数を指定しない関数の等式
は\newterm{ポイントフリー}と表現される。（ちなみに、ポイントフリー等式はしばしば関数の
合成を含み、それは点で象徴されるため、初心者には少し混乱するかもしれない。）

C++でモノイドを宣言するに最も近い方法は、C++20標準のコンセプト機能を使うことだろう。

\begin{snip}{cpp}
template<typename T>
struct mempty;

template<typename T>
T mappend(T, T) = delete;

template<typename M>
concept Monoid = requires (M m) {
    { mempty<M>::value() } -> std::same_as<M>;
    { mappend(m, m) } -> std::same_as<M>;
};
\end{snip}
最初の定義は各特殊化の単位元を保持するための構造体だ。

キーワード\code{delete}はデフォルト値が定義されていないことを意味する：個別に指定しな
ければならない。同様に、\code{mappend}にもデフォルトはない。

コンセプト\code{Monoid}は、与えられた型\code{M}に対して\code{mempty}と\code{mappend}の
適切な定義が存在するかどうかをテストする。

Monoidコンセプトのインスタンス化は、適切な特殊化とオーバーロードを提供することで達成
できる：

\begin{snip}{cpp}
template<>
struct mempty<std::string> {
    static std::string value() { return ""; }
};

template<>
std::string mappend(std::string s1, std::string s2) {
    return s1 + s2;
}
\end{snip}

\section{圏としてのモノイド}

それは集合の要素の観点からのモノイドの「親しみやすい」定義だった。しかしご存知のように、
圏論では集合とその要素から離れ、代わりに対象と射について話そうとする。だから少し観点を
変えて、二項演算の適用を集合の周りのものを「移動」または「シフト」することとして考えよう。

例えば、すべての自然数に5を加算する演算がある。それは0を5に、1を6に、2を7に、というよ
うにマップする。これは自然数の集合上で定義された関数だ。良い：関数と集合がある。一般に、
任意の数$n$に対して$n$を加算する関数――$n$の「加算器」がある。

加算器はどのように合成されるか？5を加算する関数と7を加算する関数の合成は12を加算する
関数だ。だから加算器の合成は加算の規則と等価にできる。それも良い：加算を関数合成に置き
換えられる。

しかし待て、まだある：単位元ゼロの加算器もある。ゼロを加算してもものは動かないので、
それは自然数の集合における恒等関数だ。

加算の伝統的な規則を伝える代わりに、情報を失うことなく加算器の合成の規則を伝えることも
できた。加算器の合成が結合的であることに注意しよう。なぜなら関数の合成は結合的だからだ。
そして恒等関数に対応するゼロの加算器がある。

洞察力のある読者は、整数から加算器へのマッピングが\code{mappend}の型シグネチャの2番目
の解釈\code{m -> (m -> m)}から従うことに気づいたかもしれない。これは\code{mappend}がモ
ノイド集合の要素をその集合上で作用する関数にマップすることを伝えている。

今、自然数の集合を扱っていることを忘れ、単に1つの対象、多くの射――加算器――を持つ塊と
して考えてほしい。モノイドは1対象の圏だ。実際、モノイドという名前はギリシャ語の
\emph{mono}（単一を意味する）に由来する。すべてのモノイドは合成の適切な規則に従う射の
集合を持つ1対象の圏として記述できる。

\begin{figure}[H]
  \centering
  \includegraphics[width=0.35\textwidth]{images/monoid.jpg}
\end{figure}

\noindent
文字列連結は興味深いケースだ。右のアペンダーと左のアペンダー（または前置アペンダー）を
定義する選択があるからだ。2つのモデルの合成表は互いに鏡像になっている。「foo」の後に
「bar」を追加することが、「bar」を前置した後に「foo」を前置することに対応することを容易
に確認できる。

すべての圏論的モノイド――1対象の圏――が一意な二項演算付き集合モノイドを定義するかどう
かを問うかもしれない。実は、1対象の圏から常に集合を抽出できる。この集合が射の集合だ――
我々の例では加算器。言い換えれば、圏$\cat{M}$の唯一の対象$m$のホム集合$\cat{M}(m, m)$
を持つ。この集合における二項演算を簡単に定義できる：2つの集合要素のモノイド積は対応す
る射の合成に対応する要素だ。$f$と$g$に対応する$\cat{M}(m, m)$の2つの要素を与えられた
とき、それらの積は合成$f \circ g$に対応する。これらの射の源と標的は同じ対象なので、合成
は常に存在する。そして圏の規則により結合的だ。恒等射はこの積の単位元だ。だから圏モノイ
ドから常に集合モノイドを回復できる。すべての実際の目的において、それらは同一だ。

\begin{figure}[H]
  \centering
  \includegraphics[width=0.4\textwidth]{images/monoidhomset.jpg}
  \caption{射として、および集合の点として見たモノイドのホム集合。}
\end{figure}

\noindent
数学者が指摘すべき小さな難点が1つある：射は必ずしも集合を形成する必要はない。圏の世界
には集合より大きいものがある。任意の2つの対象間の射が集合を形成する圏は局所的に小さい
と呼ばれる。約束したとおり、そのような細部はほとんど無視するが、記録のために言及してお
くべきだと思った。

圏論における多くの興味深い現象は、ホム集合の要素が合成の規則に従う射としても、集合の
点としても見られるという事実に根ざしている。ここで、$\cat{M}$における射の合成は集合
$\cat{M}(m, m)$におけるモノイド積に変換される。

\section{課題}

\begin{enumerate}
  \tightlist
  \item
        次から自由圏を生成せよ：

        \begin{enumerate}
          \tightlist
          \item
                1つのノードと辺なしのグラフ
          \item
                1つのノードと1つの（有向）辺を持つグラフ（ヒント：この辺は自分自身と
                合成できる）
          \item
                2つのノードとそれらの間の1つの矢印を持つグラフ
          \item
                1つのノードとアルファベットの文字でマークされた26個の矢印を持つグラフ：
                a, b, c \ldots{} z。
        \end{enumerate}
  \item
        これはどのような種類の順序か？

        \begin{enumerate}
          \tightlist
          \item
                包含関係を持つ集合の集合：$A$のすべての要素が$B$の要素でもある場合、$A$は
                $B$に包含される。
          \item
                次のサブタイプ関係を持つC++の型：コンパイルエラーを引き起こすことなく
                \code{T1}へのポインタを\code{T2}へのポインタを期待する関数に渡せる場合、
                \code{T1}は\code{T2}のサブタイプだ。
        \end{enumerate}
  \item
        \code{Bool}が\code{True}と\code{False}の2つの値の集合であることを考慮して、
        それが演算子\code{\&\&}（AND）と\code{||}（OR）に対してそれぞれ2つの（集合論
        的）モノイドを形成することを示せ。
  \item
        AND演算子を使った\code{Bool}モノイドを圏として表現せよ：射とその合成の規則を
        リストせよ。
  \item
        3を法とする加算をモノイド圏として表現せよ。
\end{enumerate}
